\documentclass{reportkit}

\setreportkitleftheader{REPORTKIT / ALGORITHM VISUALIZATION TEST}
\setreportkitfooter{Algorithm visualization acceptance test}
\setreportkitversion{v1.10.0}

\title{ReportKit -- algorithm and execution-state visualization primitives}
\author{Test build}
\date{17 September 2026}

\begin{document}
\maketitle

\begin{execsummary}
This is an acceptance test for reportkit-algorithm-viz.sty's Phase 1
primitives: the shared algorithm-state vocabulary, \texttt{arraystate},
\texttt{windowstate}, and \texttt{algorithmtrace}
(docs/superpowers/specs/2026-09-16-reportkit-algorithm-visualization-primitives-spec.md).
It is distinct from \texttt{algorithmblock} (control-flow pseudocode, see
algorithm\_acceptance\_test.tex): this primitive family visualizes the state
of the data \emph{while} an algorithm executes.
\end{execsummary}

\section{Two-pointer array scan}
The motivating \texttt{arraystate} example: indices, left/right cursors, and
an active-range highlight.

\begin{diagram}[type=array,width=\textwidth,caption={Two-pointer elimination narrows the active range from both ends.},description={A six-element array with a left pointer at index two and a right pointer at index five; the range between them is marked active.}]
\begin{arraystate}
  \cell{-4}
  \cell{-1}
  \cell{-1}
  \cell{0}
  \cell{1}
  \cell{2}
  \pointer[below]{L}{2}
  \pointer[below]{R}{5}
  \range[state=active]{2}{5}
\end{arraystate}
\end{diagram}

\section{Sliding window}
\texttt{windowstate}, a thin composition over \texttt{arraystate}: the same
active-range highlight, automatic left/right cursors, and entering/leaving
annotations.

\begin{diagram}[type=array,width=\textwidth,caption={A three-element sliding window advances by one position.},description={A six-element array with a window spanning indices two to four, an incoming value of five, and an outgoing value of one.}]
\begin{windowstate}
  \values{4,2,7,1,3,6}
  \window{2}{4}
  \entering{5}
  \leaving{1}
\end{windowstate}
\end{diagram}

\section{Prefix sums: multi-row alignment}
Multi-row \texttt{arraystate} via \texttt{\string\row}, with a range on the
named \texttt{values} row and a running-total annotation.

\begin{diagram}[type=array,width=\textwidth,caption={A prefix-sum row aligned under its source values.},description={A four-element values row aligned above a five-element prefix-sum row, with indices one to two of the values row marked active.}]
\begin{arraystate}
  \row{values}{3,1,4,2}
  \row{prefix}{0,3,4,8,10}
  \range[row=values,state=active]{1}{2}
  \annotation{$P_3 - P_1 = 5$}
\end{arraystate}
\end{diagram}

\section{Full state vocabulary}
All nine shared algorithm states in one row, so state meaning stays
distinguishable through border weight and line style, not fill color alone.

\begin{diagram}[type=array,width=\textwidth,caption={Every shared algorithm state applied to one cell each.},description={Nine cells, one per shared algorithm state: current, active, candidate, frontier, visited, resolved, discarded, blocked, and unseen.}]
\begin{arraystate}[indices=false]
  \cell[state=current]{cur}
  \cell[state=active]{act}
  \cell[state=candidate]{cand}
  \cell[state=frontier]{front}
  \cell[state=visited]{vis}
  \cell[state=resolved]{res}
  \cell[state=discarded]{disc}
  \cell[state=blocked]{blk}
  \cell[state=unseen]{uns}
\end{arraystate}
\end{diagram}

\section{Algorithmtrace: ordered snapshots}
Three ordered snapshots of a sliding-window advance, read left to right.

\begin{diagram}[type=trace,width=\textwidth,caption={A three-element window advances across a six-element array.},description={Three ordered snapshots show a three-element window at positions zero to two, one to three, and two to four.}]
\begin{algorithmtrace}[columns=3]
  \snapshot{Initial}{%
    \begin{arraystate}[cell width=.85,row height=1.0]
      \cell{4}\cell{2}\cell{7}\cell{1}\cell{3}\cell{6}
      \range[state=active]{0}{2}
    \end{arraystate}%
  }
  \snapshot{Advance right}{%
    \begin{arraystate}[cell width=.85,row height=1.0]
      \cell{4}\cell{2}\cell{7}\cell{1}\cell{3}\cell{6}
      \range[state=active]{1}{3}
    \end{arraystate}%
  }
  \snapshot{Advance again}{%
    \begin{arraystate}[cell width=.85,row height=1.0]
      \cell{4}\cell{2}\cell{7}\cell{1}\cell{3}\cell{6}
      \range[state=active]{2}{4}
    \end{arraystate}%
  }
\end{algorithmtrace}
\end{diagram}

\section{Grayscale and semantic check}
\begin{principle}{State is not just color}
Every shared algorithm state carries a distinct line-style fragment (bold
solid for \texttt{current}, dashed for \texttt{candidate} and
\texttt{discarded}, dotted for \texttt{blocked}) in addition to its
draw/fill/text colors, so the diagram above should stay legible if printed
in grayscale.
\end{principle}

\end{document}
